\documentclass[11pt]{article}

\usepackage[margin=1in]{geometry}
\usepackage{booktabs}
\usepackage{array}
\usepackage{hyperref}

\title{Avaliacao do GPT-OSS-20B com Chain-of-Thought e Process Supervision Promptado}
\author{Evaluation harness local com vLLM}
\date{Rodadas principais: 20260510-185218 e 20260510-231913}

\begin{document}
\maketitle

\begin{abstract}
Este relatorio consolida duas avaliacoes locais do modelo
\texttt{openai/gpt-oss-20b}. A primeira compara prompting few-shot padrao
contra few-shot com \emph{chain-of-thought} (CoT), inspirada em
\emph{Chain-of-Thought Prompting Elicits Reasoning in Large Language Models}.
A segunda implementa uma aproximacao promptada de \emph{Process Supervision},
inspirada em \emph{Let's Verify Step by Step}, na qual o mesmo modelo gera
solucoes para problemas MATH e tambem julga cada passo via logprobs. Todos os
experimentos foram executados localmente com vLLM em quatro GPUs RTX 4000 Ada.
Os resultados indicam que CoT melhora GSM8K de 94,5\% para 96,5\%, mas nao
ajuda em tarefas simbolicas simples. No experimento de Process Supervision, o
best-of-8 melhora a acuracia de 86,0\% para 96,0\%, mas o PRM promptado empata
com majority vote e nao supera o baseline simples.
\end{abstract}

\section{Configuracao experimental}

As avaliacoes usaram o modelo \texttt{openai/gpt-oss-20b} servido localmente
por vLLM com API OpenAI-compatible. O servidor foi executado em quatro GPUs
NVIDIA RTX 4000 Ada Generation com tensor parallel igual a 4. O objetivo nao
foi reproduzir todos os numeros dos papers originais, mas implementar
experimentos locais controlados que permitissem estudar dois mecanismos:
prompting com CoT e selecao por supervisao de processo.

\begin{table}[h]
\centering
\caption{Configuracao geral dos experimentos.}
\label{tab:general_setup}
\begin{tabular}{ll}
\toprule
Campo & Valor \\
\midrule
Modelo & \texttt{openai/gpt-oss-20b} \\
Backend & vLLM OpenAI-compatible local \\
Paralelismo & Tensor parallel em 4 GPUs \\
GPUs & 4x NVIDIA RTX 4000 Ada Generation \\
Hardware alvo & Inferencia local distribuida \\
Raw CoT interno & Nao utilizado no relatorio \\
Artefatos CoT & \texttt{experiments/01\_chain\_of\_thought/artifacts/20260510-185218/} \\
Artefatos PRM & \texttt{experiments/02\_process\_supervision/artifacts/20260510-231913/} \\
\bottomrule
\end{tabular}
\end{table}

Em ambos os estudos, a avaliacao se baseia apenas em respostas visiveis,
respostas parseadas e metricas agregadas. O relatorio nao usa nem expande raw
chain-of-thought interno do modelo.

\section{Experimento 1: CoT versus resposta direta}

\subsection{Metodo}

O primeiro experimento compara duas condicoes de prompting. A condicao
\emph{Sem CoT} usa exemplos few-shot com resposta direta. A condicao
\emph{CoT} usa exemplos few-shot com raciocinio intermediario antes da resposta
final. As duas condicoes compartilham o mesmo modelo, backend, temperatura
zero, \texttt{reasoning\_effort=medium} e o mesmo contrato de saida:
\texttt{Final answer: <answer>}. Esse contrato reduz a chance de o parser
extrair numeros intermediarios em vez da resposta final.

Foram avaliadas tres tarefas: GSM8K, Last Letter e Coin Flip. GSM8K representa
raciocinio aritmetico multi-etapa; Last Letter e uma tarefa simbolica local; e
Coin Flip avalia rastreamento simples de estado binario.

GSM8K e um benchmark de problemas matematicos em linguagem natural criado pela
OpenAI para avaliar raciocinio aritmetico multi-etapa em modelos de linguagem.
O nome vem de ``Grade School Math 8K'', pois o conjunto contem cerca de 8,5 mil
problemas no estilo matematica de ensino fundamental, normalmente resolviveis
com operacoes como soma, subtracao, multiplicacao, divisao, proporcoes e
contagem. Cada exemplo traz um enunciado textual e uma solucao com resposta
final numerica. Um problema tipico seria: ``Joao tinha 12 macas, comprou mais 8
e depois dividiu igualmente entre 5 amigos; com quantas macas cada amigo
ficou?'', exigindo que o modelo decomponha o texto em etapas aritmeticas antes
de chegar a resposta. No nosso experimento, usamos o dataset publico
\texttt{openai/gsm8k}, disponivel no Hugging Face, especificamente a
configuracao \texttt{main} e o split \texttt{test}. Do conjunto completo de
aproximadamente 8,5 mil problemas, avaliamos uma amostra de 200 problemas do
split de teste. A resposta correta foi extraida do campo de solucao do dataset,
onde o valor final aparece apos o marcador \texttt{\#\#\#\#}, e o modelo foi
considerado correto quando sua linha \texttt{Final answer: <answer>} correspondia
numericamente a esse valor.

A tarefa Last Letter foi usada como uma avaliacao simbolica simples, inspirada
nas tarefas sinteticas utilizadas no artigo de Chain-of-Thought. Nessa tarefa,
o modelo recebe uma sequencia de palavras ou nomes e deve retornar a
concatenacao da ultima letra de cada termo, preservando a ordem em que os termos
aparecem. Por exemplo, para a entrada ``Donald Hamilton'', a resposta esperada
e \texttt{dn}, pois a ultima letra de ``Donald'' e \texttt{d} e a ultima letra
de ``Hamilton'' e \texttt{n}. Diferentemente do GSM8K, essa tarefa nao vem de
um dataset externo fixo: ela foi gerada localmente pelo nosso harness, de forma
deterministica a partir de uma seed, para permitir controle total sobre os
exemplos e as respostas ouro. No experimento, avaliamos 100 exemplos de Last
Letter. A resposta do modelo foi extraida da linha \texttt{Final answer:
<answer>} e comparada exatamente, apos normalizacao textual, com a string ouro
gerada pela regra da tarefa.

A tarefa Coin Flip tambem foi gerada localmente e segue o estilo de tarefas
simbolicas de rastreamento de estado usadas para avaliar Chain-of-Thought. Nessa
tarefa, o enunciado descreve o estado inicial de uma moeda e uma sequencia de
acoes, como virar ou nao virar a moeda. O modelo deve determinar se, ao final da
sequencia, a moeda esta com a face inicial para cima, respondendo em formato
\texttt{yes} ou \texttt{no}. Um exemplo simplificado seria: ``A coin is heads
up. Ana flips the coin. Ana flips the coin again. Is the coin still heads
up?''; como duas viradas retornam a moeda ao estado inicial, a resposta correta
seria \texttt{yes}. No nosso experimento, avaliamos 100 exemplos de Coin Flip,
tambem gerados deterministicamente por seed. A resposta ouro foi calculada por
regra, simulando as mudancas de estado da moeda, e a saida do modelo foi
avaliada por um parser de respostas \texttt{yes}/\texttt{no}.



\begin{table}[h]
\centering
\caption{Tarefas usadas no experimento CoT.}
\label{tab:cot_tasks}
\begin{tabular}{p{0.16\linewidth}p{0.22\linewidth}p{0.52\linewidth}}
\toprule
Tarefa & Origem & Descricao \\
\midrule
GSM8K & \texttt{openai/gsm8k}, split \texttt{test}, config \texttt{main} &
Conjunto de problemas matematicos de ensino fundamental. Cada item contem um
enunciado textual e uma resposta numerica final. A avaliacao usa 200 exemplos
e extrai o \textit{gold} do campo de resposta apos o marcador \texttt{\#\#\#\#}. \\
Last Letter & Gerada localmente, inspirada no paper de CoT &
Tarefa simbolica em que o modelo recebe uma lista de palavras ou nomes e deve
concatenar a ultima letra de cada termo. Por exemplo, para dois nomes, a saida
esperada e uma string com duas letras. A rodada usa 100 exemplos gerados de
forma deterministica por seed. \\
Coin Flip & Gerada localmente, inspirada no paper de CoT &
Tarefa de rastreamento de estado. O enunciado descreve uma moeda inicialmente
em um estado e uma sequencia de acoes de virar ou nao virar a moeda. O modelo
deve responder se a moeda termina com a face inicial para cima, em formato
\texttt{yes}/\texttt{no}. A rodada usa 100 exemplos gerados de forma
deterministica por seed. \\
\bottomrule
\end{tabular}
\end{table} 
\newpage
\subsection{Resultados}

\begin{table}[h]
\centering
\caption{Acuracia por tarefa e condicao no experimento CoT.}
\label{tab:cot_accuracy}
\begin{tabular}{lrrrr}
\toprule
Tarefa & Condicao & Corretas/N & Acuracia & Delta CoT \\
\midrule
GSM8K & Sem CoT & 189/200 & 94,5\% & -- \\
GSM8K & CoT & 193/200 & 96,5\% & +2,0 p.p. \\
Last Letter & Sem CoT & 85/100 & 85,0\% & -- \\
Last Letter & CoT & 84/100 & 84,0\% & -1,0 p.p. \\
Coin Flip & Sem CoT & 100/100 & 100,0\% & -- \\
Coin Flip & CoT & 100/100 & 100,0\% & +0,0 p.p. \\
\bottomrule
\end{tabular}
\end{table}

\begin{table}[h]
\centering
\caption{Custo medio por item no experimento CoT.}
\label{tab:cot_cost}
\begin{tabular}{lrrrr}
\toprule
Tarefa/condicao & Latencia (s) & Prompt tok. & Resposta tok. & Total tok. \\
\midrule
GSM8K Sem CoT & 1,03 & 584,4 & 202,1 & 786,5 \\
GSM8K CoT & 1,12 & 892,4 & 219,0 & 1111,4 \\
Last Letter Sem CoT & 0,38 & 303,4 & 72,8 & 376,1 \\
Last Letter CoT & 0,75 & 428,4 & 146,7 & 575,1 \\
Coin Flip Sem CoT & 1,51 & 502,3 & 298,2 & 800,5 \\
Coin Flip CoT & 0,82 & 852,3 & 158,4 & 1010,8 \\
\bottomrule
\end{tabular}
\end{table}

\subsection{Analise critica}

O principal ganho de CoT aparece em GSM8K: a acuracia sobe de 94,5\% para
96,5\%. A comparacao pareada mostrou que CoT corrigiu seis itens que a resposta
direta errou, enquanto perdeu dois itens que a resposta direta acertou. Esse
padrao e consistente com a hipotese de que demonstracoes intermediarias ajudam
em problemas aritmeticos que exigem decomposicao. Ainda assim, o ganho e
pequeno e a amostra nao e suficiente para afirmar uma vantagem universal.

Em Last Letter, CoT nao melhorou o desempenho. A tarefa exige uma operacao
simbolica curta e exata; nesse caso, raciocinio intermediario pode apenas criar
mais pontos de erro. Em Coin Flip, ambas as condicoes atingiram 100\%, o que
torna a tarefa pouco informativa para distinguir os metodos nesta configuracao.

Do ponto de vista de custo, CoT aumenta o numero medio de tokens. Em GSM8K, o
total medio cresce de 786,5 para 1111,4 tokens por item, aumento aproximado de
41\%. Assim, o uso de CoT parece justificado para raciocinio matematico
multi-etapa, mas nao para tarefas simbolicas simples ou ja saturadas.

\section{Experimento 2: \textit{Process Supervision} promptado}

\subsection{Metodo}

O segundo experimento implementa uma aproximacao promptada de Process
Supervision. O GPT-OSS-20B gera oito solucoes por problema MATH. Em seguida, o
mesmo modelo julga cada passo de cada solucao com os labels
\texttt{positive}, \texttt{neutral} e \texttt{negative}. As probabilidades sao
estimadas com logprobs do endpoint de completions. Seguindo a configuracao
principal descrita no paper de referencia, \texttt{neutral} e tratado como
positivo, e o score da solucao e o produto dos scores de seus passos.

Este experimento nao treina um PRM real e nao usa labels humanos de PRM800K.
Ele mede a utilidade operacional de um verificador promptado com o mesmo modelo
que gerou as solucoes. Isso cria risco de erros correlacionados: uma solucao
errada, mas plausivel para o gerador, tambem pode parecer plausivel para o
verificador.

O benchmark usado foi o MATH, tambem conhecido nesta implementacao pelo dataset
\texttt{EleutherAI/hendrycks\_math} no Hugging Face. Esse conjunto foi criado
para avaliar resolucao de problemas matematicos mais dificeis que GSM8K, em
estilo olimpiada, competicao escolar e cursos avancados de ensino medio. Os
enunciados cobrem areas como algebra, geometria, teoria dos numeros,
precalculo, contagem e probabilidade, e normalmente exigem uma sequencia de
passos matematicos, nao apenas uma operacao aritmetica direta. Um exemplo
tipico de problema MATH pode pedir, por exemplo, a simplificacao de uma
expressao algebrica, a area de uma figura geometrica ou a solucao de uma
equacao com restricoes. No nosso experimento, usamos apenas o split
\texttt{test} e quatro configuracoes do dataset:
\texttt{intermediate\_algebra}, \texttt{precalculus}, \texttt{geometry} e
\texttt{number\_theory}. Esses subconjuntos somaram 2468 itens vistos pelo
loader. Como a avaliacao automatica dependia de comparacao numerica, filtramos
os exemplos para manter apenas aqueles em que a resposta final numerica podia
ser extraida de uma expressao \verb|\boxed{...}|; nessa etapa, 781 itens foram
descartados. A partir dos itens restantes, selecionamos 50 problemas usando a
seed 43. Para cada problema, o modelo gerou 8 solucoes candidatas, totalizando
400 solucoes avaliadas pelo seletor PRM, majority vote e oracle.

Antes da rodada final, o experimento passou por uma etapa incremental de
implementacao e validacao. Primeiro, foi implementado um pipeline minimo para
gerar solucoes, quebrar cada solucao em passos, avaliar cada passo com um
verificador promptado e salvar os artefatos em \texttt{JSONL}, \texttt{CSV}, e
\texttt{Markdown}. Em seguida, rodamos testes menores, incluindo uma
rodada inicial com poucos problemas e uma rodada \texttt{10x4}, para validar o
formato das respostas, a escrita dos arquivos, a agregacao das metricas e o
comportamento dos seletores \emph{first}, \emph{majority vote},
\emph{PRM best-of-N} e \emph{oracle best-of-N}. Essas rodadas preliminares
tambem revelaram problemas praticos importantes, especialmente truncamento de
respostas e erros de parsing de fracoes finais. Por isso, antes da rodada
\texttt{50x8}, aumentamos o orcamento de geracao para
\texttt{max\_tokens=4096}, ajustamos a temperatura para 0,9 a fim de aumentar a
diversidade dos candidatos, e corrigimos o parser para reconhecer respostas
como \texttt{1/3}, \verb|\frac{5}{9}| e expressoes com \verb|\displaystyle|.
Assim, os resultados reportados representam a rodada final apos a estabilizacao
do pipeline, nao apenas uma execucao unica isolada.


\begin{table}[h]
\centering
\caption{Configuracao da rodada PRM.}
\label{tab:prm_setup}
\begin{tabular}{ll}
\toprule
Campo & Valor \\
\midrule
Benchmark & \texttt{EleutherAI/hendrycks\_math} \\
Configs MATH & intermediate algebra, precalculus, geometry, number theory \\
Filtro & Apenas respostas numericas extraiveis \\
Problemas & 50 \\
Solucoes por problema & 8 \\
Candidatos totais & 400 \\
Temperatura de geracao & 0,9 \\
\texttt{max\_tokens} & 4096 \\
Seed & 43 \\
Scoring PRM & Produto dos scores por passo \\
\bottomrule
\end{tabular}
\end{table}

O protocolo de selecao foi inspirado no desenho experimental de
\emph{Let's Verify Step by Step}~\cite{lightman2023verify}, que avalia nao
apenas uma unica resposta do modelo, mas tambem estrategias para selecionar a
melhor solucao entre varias amostras. Por isso, comparamos quatro metodos. O
primeiro, \emph{first sample}, usa simplesmente a primeira solucao gerada e
serve como baseline sem busca. O segundo, \emph{majority vote}, escolhe a
resposta final mais frequente entre as oito solucoes candidatas; esse baseline
é importante porque muitos ganhos de \emph{best-of-N} podem vir apenas da
agregacao de varias amostras, sem qualquer verificador por passo. O terceiro,
\emph{PRM best-of-N}, seleciona a solucao com maior score de processo, isto e,
a solucao cujos passos foram julgados como mais confiaveis pelo verificador
promptado. Por fim, \emph{oracle best-of-N} escolhe uma solucao correta sempre
que ela existe entre as candidatas; esse metodo nao e utilizavel na pratica,
mas funciona como teto superior para medir quanta margem de melhoria ainda
existiria se o seletor fosse perfeito.


\subsection{Resultados}

\begin{table}[h]
\centering
\caption{Acuracia corrigida por metodo de selecao no experimento PRM.}
\label{tab:prm_accuracy}
\begin{tabular}{lrrrr}
\toprule
Metodo & Corretas/N & Acuracia & Falhas parse & Cobertura parse \\
\midrule
First sample & 43/50 & 86,0\% & 5 & 90,0\% \\
Majority vote & 48/50 & 96,0\% & 0 & 100,0\% \\
PRM best-of-N & 48/50 & 96,0\% & 0 & 100,0\% \\
Oracle best-of-N & 49/50 & 98,0\% & 1 & 98,0\% \\
\bottomrule
\end{tabular}
\end{table}

\begin{table}[h]
\centering
\caption{Custo medio por problema selecionado no experimento PRM.}
\label{tab:prm_cost}
\begin{tabular}{lrrrr}
\toprule
Metodo & Lat. geracao (s) & Lat. PRM (s) & Tok. geracao & Tok. PRM \\
\midrule
First sample & 7,94 & 0,00 & 1794,4 & 0,0 \\
Majority vote & 61,67 & 0,00 & 14017,2 & 0,0 \\
PRM best-of-N & 61,67 & 2,14 & 14017,2 & 39707,4 \\
Oracle best-of-N & 61,67 & 0,00 & 14017,2 & 0,0 \\
\bottomrule
\end{tabular}
\end{table}

\begin{table}[h]
\centering
\caption{Diagnostico dos candidatos gerados no experimento PRM.}
\label{tab:prm_diagnostic}
\begin{tabular}{lr}
\toprule
Metrica & Valor \\
\midrule
Candidatos totais & 400 \\
Candidatos corretos apos correcao do parser & 345 \\
Falhas de parse em candidatos & 39 \\
Problemas com 8/8 candidatos corretos & 36 \\
Problemas com mistura de candidatos corretos e incorretos & 13 \\
Problemas sem nenhum candidato correto & 1 \\
Problemas com uma unica resposta parseavel distinta & 43 \\
AUC do score PRM para separar candidato correto/incorreto & 0,397 \\
\bottomrule
\end{tabular}
\end{table}

\subsection{Analise critica}

O resultado mais importante e que gerar oito solucoes por problema melhora
substancialmente o desempenho: o primeiro candidato acerta 86,0\%, enquanto
majority vote acerta 96,0\%. Isso mostra que o modelo frequentemente contem a
resposta correta no conjunto de candidatos mesmo quando a primeira amostra falha.

Entretanto, o PRM promptado nao supera majority vote. Apos a correcao do
parser, os dois metodos empatam em 48/50. O oracle chega a 49/50, mostrando que
havia uma oportunidade adicional de selecao que nem majority vote nem PRM
capturaram. O caso principal foi um problema de precalculo em que havia apenas
um candidato correto entre oito; tanto a votacao quanto o PRM escolheram uma
resposta incorreta. Esse caso ilustra a limitacao do verificador promptado:
quando a solucao correta e rara e longa, o score por produto pode penaliza-la.

O diagnostico de scores reforca essa cautela. Embora o score medio de
candidatos corretos seja maior que o de candidatos incorretos em alguns
resumos, a ordenacao fina nao e confiavel: a AUC do score PRM para separar
candidatos corretos de incorretos ficou em aproximadamente 0,397. Isso indica
que o score bruto nao deve ser interpretado como probabilidade calibrada de
correcao. Uma causa provavel e o uso do produto dos scores por passo, que tende
a favorecer solucoes mais curtas. De fato, o PRM selecionou solucoes com media
de 8,86 passos, enquanto majority vote selecionou solucoes com media de 14,06
passos.

O custo operacional do PRM promptado tambem e relevante. Majority vote exige
apenas os tokens de geracao dos oito candidatos. PRM best-of-N exige os mesmos
tokens de geracao e ainda adiciona cerca de 39707 tokens de verificacao por
problema. Portanto, nesta configuracao, o PRM promptado aumenta muito o custo
sem melhorar a acuracia sobre majority vote.

\section{Discussao geral}

Os dois experimentos devem ser lidos como uma demonstracao didatica da evolucao
recente em modelos de linguagem: primeiro, observa-se que induzir o modelo a
explicitar passos intermediarios por meio de \emph{chain-of-thought} pode
melhorar respostas em tarefas que exigem raciocinio; em seguida, esses passos
intermediarios passam a ser tratados nao apenas como texto auxiliar, mas como
objetos que podem ser avaliados, comparados e usados como sinal de treinamento.
No Experimento 1, a melhoria em GSM8K mostra o fenomeno original do CoT: mesmo
sem alterar os pesos do modelo, mudar o formato do prompt desloca o
comportamento do modelo para solucoes mais decompostas e frequentemente mais
corretas. Esse resultado e especialmente claro em problemas aritmeticos
multi-etapa, nos quais a resposta direta pode falhar por pular uma etapa de
calculo ou interpretar mal uma relacao do enunciado.

O Experimento 2 representa o passo conceitual seguinte. A partir do momento em
que uma solucao passa a ser escrita como uma sequencia de passos, torna-se
possivel avaliar o processo, e nao apenas a resposta final. Essa e a motivacao
central de Process Supervision: transformar o raciocinio intermediario em um
sinal de supervisao para selecionar melhores solucoes e, em sistemas maiores,
alimentar etapas de post-training e pipelines de aprendizado por reforco. Nesta
implementacao, nao treinamos um PRM real; usamos um PRM promptado para mostrar
operacionalmente como esse sinal poderia ser produzido. Mesmo assim, o
experimento ilustra a arquitetura da ideia: gerar varias solucoes, decompor em
passos, julgar cada passo, agregar os scores e escolher uma solucao por
best-of-N.

Os resultados tambem mostram por que esse tipo de supervisao se tornou
importante. No experimento PRM, o conjunto de oito candidatos por problema
frequentemente continha uma resposta correta mesmo quando a primeira amostra
falhava. Isso evidencia que parte do problema nao e apenas ``o modelo sabe ou
nao sabe'', mas tambem como selecionar, entre varias trajetorias possiveis, a
trajetoria mais confiavel. O majority vote ja foi um baseline forte nesta
rodada, e o PRM promptado nao o superou; ainda assim, essa limitacao e
informativa. Ela mostra que simplesmente pedir ao mesmo modelo para julgar seus
proprios passos nao equivale a ter um reward model treinado com supervisao
humana. Em outras palavras, o experimento reproduz em pequena escala a
motivacao para PRMs treinados: e necessario transformar avaliacoes de processo
em um sinal mais calibrado, robusto e util para selecao e treinamento.

Assim, o valor principal do projeto nao esta apenas nas acuracias finais, mas
na cadeia experimental que elas tornam visivel. CoT demonstra que a forma de
externalizar o raciocinio altera o desempenho do modelo. Process Supervision
mostra como esse raciocinio externalizado pode ser convertido em dados de
avaliacao por passo. Essa passagem, de ``promptar para raciocinar'' para
``supervisionar o processo de raciocinio'', e uma das bases conceituais que
influenciaram novas formas de post-training, reward modeling e aprendizado por
reforco em modelos de linguagem.

\section{Limitacoes}

Este estudo tem quatro limitacoes principais. Primeiro, as amostras sao
exploratorias: 200 itens de GSM8K, 100 itens por tarefa simbolica e 50 problemas
MATH no PRM. Segundo, o PRM e promptado, nao treinado com labels humanos; logo,
nao deve ser confundido com um reward model supervisionado real. Terceiro, o
mesmo modelo gera e julga, aumentando o risco de erros correlacionados. Quarto,
o filtro numeric-only simplifica MATH e remove exemplos cujo formato de resposta
nao era facilmente comparavel por regra.

\section{Conclusao}

Para o GPT-OSS-20B executado localmente, CoT mostrou beneficio moderado em
GSM8K, mas nao trouxe ganho em tarefas simbolicas simples. No estudo de Process
Supervision, o best-of-8 foi efetivo, mas o PRM promptado nao superou majority
vote e adicionou custo substancial. Assim, os resultados sustentam uma
recomendacao pragmatica: usar CoT seletivamente para problemas matematicos
multi-etapa e usar best-of-N com majority vote como baseline forte antes de
adotar verificadores por passo. Para demonstrar ganho real de Process
Supervision, o proximo passo deveria comparar agregadores alternativos de score
por passo, como media de log-score ou minimo por passo, ou treinar um PRM real
com labels humanos.

\begin{thebibliography}{9}

\bibitem{wei2022cot}
Jason Wei et al.
\newblock Chain-of-Thought Prompting Elicits Reasoning in Large Language Models.
\newblock arXiv:2201.11903, 2022.

\bibitem{lightman2023verify}
Hunter Lightman et al.
\newblock Let's Verify Step by Step.
\newblock arXiv:2305.20050, 2023.

\bibitem{gsm8k}
OpenAI.
\newblock GSM8K dataset.
\newblock \url{https://huggingface.co/datasets/openai/gsm8k}.

\bibitem{math}
EleutherAI.
\newblock Hendrycks MATH dataset.
\newblock \url{https://huggingface.co/datasets/EleutherAI/hendrycks_math}.

\end{thebibliography}

\end{document}
